\chapter*{Inleiding}

Deze studies richten zich op het probleem van de hermeneutiek. Het fenomeen van het begrijpen en van de juiste interpretatie van wat begrepen is, is geen probleem dat zich beperkt tot de methodologie van de geesteswetenschappen alleen. Er bestond al lang een theologische en een juridische hermeneutiek, die niet zozeer theoretisch waren, maar veeleer een bijproduct en een hulpmiddel vormden voor de praktische activiteit van de rechter of de geestelijke die zijn theoretische opleiding had afgerond. Al vanaf haar historische oorsprong gaat het probleem van de hermeneutiek verder dan de grenzen van het methodenbegrip zoals dat door de moderne wetenschap is gesteld. Het begrijpen en interpreteren van teksten is niet louter een zaak van de wetenschap, maar hoort onmiskenbaar bij de menselijke ervaring van de wereld in het algemeen. Het hermeneutische fenomeen is in wezen helemaal geen methodisch probleem. Het gaat niet om een methode van begrip waarlangs teksten, net als alle andere ervaringsobjecten, aan een wetenschappelijke analyse worden onderworpen. Het draait ook niet primair om het verzamelen van geverifieerde kennis, zoals die welke voldoet aan het methodologische ideaal van de wetenschap — en toch gaat het ook om kennis en om waarheid. Bij het begrijpen van traditie worden niet alleen teksten begrepen, maar worden ook inzichten verworven en waarheden gekend. Maar wat voor kennis en wat voor waarheid?

Gelet op de dominantie van de moderne wetenschap in de filosofische verheldering en rechtvaardiging van het kennisbegrip en het waarheidsbegrip, lijkt deze vraag niet legitiem. Toch is ze onvermijdelijk, zelfs binnen de wetenschappen zelf. Het fenomeen van het begrijpen doordringt niet alleen alle menselijke betrokkenheid bij de wereld; het heeft ook een eigen geldigheid binnen de wetenschap en verzet zich tegen elke poging om het in termen van wetenschappelijke methoden te herinterpreteren. De volgende onderzoeken vertrekken vanuit het verzet binnen de moderne wetenschap zelf tegen de universele claim van de wetenschappelijke methode. Ze streven ernaar de ervaring van waarheid te zoeken die voorbij het domein van de wetenschappelijke methode gaat, waar die ervaring ook te vinden is, en om haar legitimiteit te onderzoeken. Vandaar dat de geesteswetenschappen verbonden zijn met ervaringswijzen die buiten de wetenschap liggen: met de ervaringen van de filosofie, van de kunst en van de geschiedenis zelf. Dit zijn allemaal ervaringswijzen waarin een waarheid wordt meegedeeld die niet kan worden geverifieerd met de methodologische middelen die eigen zijn aan de wetenschap.

De hedendaagse filosofie is zich hier terdege van bewust. Maar het is een heel andere vraag in hoeverre de waarheidsclaim van dergelijke ervaringswijzen buiten de wetenschap filosofisch kan worden gelegitimeerd. Het huidige belang voor het hermeneutische fenomeen berust, naar ik meen, op het feit dat alleen een diepgaander onderzoek van het fenomeen van het begrijpen deze legitimatie kan bieden. Deze overtuiging wordt krachtig ondersteund door het belang dat de hedendaagse filosofie hecht aan de geschiedenis van de filosofie. Wat de historische traditie van de filosofie betreft, komt het begrijpen ons voor als een superieure ervaring die ons in staat stelt om gemakkelijk door de illusie van de historische methode heen te prikken, kenmerkend voor het onderzoek in de geschiedenis van de filosofie. Het hoort bij de elementaire ervaring van de filosofie dat wanneer we proberen de klassiekers van het filosofisch denken te begrijpen, deze vanzelf een claim op waarheid maken die het bewustzijn van latere tijden noch kan afwijzen noch kan overstijgen. Het naïeve zelfvertrouwen van het heden kan zich verzetten tegen het idee dat het filosofisch bewustzijn de mogelijkheid toegeeft dat het eigen filosofische inzicht inferieur zou kunnen zijn aan dat van Plato of Aristoteles, Leibniz, Kant of Hegel. Men zou kunnen denken dat het een zwakte is dat de hedendaagse filosofie haar klassieke erfenis probeert te interpreteren en te verwerken met deze erkenning van haar eigen zwakte. Maar het is ongetwijfeld een veel grotere zwakte voor het filosofisch denken om een dergelijk zelfonderzoek uit de weg te gaan en in plaats daarvan de Faust te willen uithangen. Het is duidelijk dat in het begrijpen van de teksten van deze grote denkers een waarheid wordt gekend die op geen andere manier bereikt had kunnen worden, ook al strookt dit tegen de maatstaf van onderzoek en vooruitgang waarlangs de wetenschap zich meet.

Hetzelfde geldt voor de ervaring van kunst. Hier is de wetenschappelijke navorsing die door de ‘wetenschap van de kunst’ wordt bedreven, zich vanaf het begin bewust dat zij noch de ervaring van kunst kan vervangen noch kan overstijgen. Het feit dat door een kunstwerk een waarheid wordt ervaren die we op geen andere manier kunnen bereiken, vormt de filosofische betekenis van kunst, die zich assertief toont tegenover alle pogingen om haar rationeel weg te redeneren. Daardoor is, samen met de ervaring van de filosofie, de ervaring van kunst de meest indringende vermaning aan het wetenschappelijk bewustzijn om haar eigen grenzen te erkennen.

Vandaar dat het volgende onderzoek begint met een kritiek op het esthetische bewustzijn, om de ervaring van waarheid die tot ons komt via het kunstwerk te verdedigen tegen de esthetische theorie die zich laat beperken tot een wetenschappelijk waarheidsbegrip. Maar dit boek stelt zich niet tevreden met het rechtvaardigen van de waarheid van kunst; in plaats daarvan probeert het, uitgaande van dit startpunt, een kennis- en waarheidsbegrip te ontwikkelen dat aansluit bij de totaliteit van onze hermeneutische ervaring. Net zoals in de ervaring van kunst het gaat om waarheden die wezenlijk verder reiken dan het bereik van methodische kennis, geldt hetzelfde voor de geesteswetenschappen in hun geheel: in hen wordt onze historische traditie in al haar vormen zeker tot object van onderzoek gemaakt, maar tegelijkertijd komt er in die traditie een waarheid tot uiting. Fundamenteel reikt de ervaring van historische traditie veel verder dan die aspecten ervan die objectief kunnen worden onderzocht. Ze is niet alleen waar of onwaar in de zin waarin de historische kritiek beslist, maar draagt altijd een waarheid in zich mee waaraan men moet proberen deel te hebben.

Daarom proberen deze studies over hermeneutiek, die vertrekken vanuit de ervaring van kunst en historische traditie, het hermeneutische fenomeen in zijn volle omvang te presenteren. Het gaat erom daarin een ervaring van waarheid te herkennen die niet alleen filosofisch gerechtvaardigd moet worden, maar die zelf een manier van filosofie beoefenen is. De hier ontwikkelde hermeneutiek is dus geen methodologie van de geesteswetenschappen, maar een poging om te begrijpen wat de geesteswetenschappen werkelijk zijn, voorbij hun methodologische zelfbewustzijn, en wat hen verbindt met de totaliteit van onze ervaring van de wereld.

Als we het begrijpen tot object van onze reflectie maken, is het doel niet een kunst of techniek van het begrijpen, zoals de traditionele literair-theologische hermeneutiek streefde te zijn. Een dergelijke kunst of techniek zou er niet in slagen te erkennen dat, gelet op de waarheid die tot ons spreekt uit de traditie, een formele techniek zich ten onrechte een superieure positie zou toemeten. Hoewel ik in het volgende zal aantonen hoeveel er van gebeurtenis werkzaam is in al het begrijpen, en hoe weinig de tradities waarin we staan verzwakt zijn door het moderne historische bewustzijn, is het niet mijn bedoeling om voorschriften op te stellen voor de wetenschappen of voor de levenspraktijk, maar om verkeerd denken over wat zij zijn te corrigeren.

Ik hoop op deze manier een inzicht te versterken dat in onze snel veranderende tijd dreigt te worden vergeten. Dingen die veranderen, dringen zich veel sterker aan onze aandacht op dan diegene die onveranderd blijven. Dat is een algemene wet van ons intellectuele leven. Daardoor lopen de perspectieven die voortvloeien uit de ervaring van historische verandering altijd het risico om overdreven te worden, omdat ze vergeten wat onzichtbaar blijft bestaan. In het moderne leven wordt ons historische bewustzijn voortdurend overprikkeld. Als gevolg daarvan — hoewel het, zoals ik hoop aan te tonen, een schadelijke kortsluiting is — reageren sommigen op deze overschatting van historische verandering door een beroep te doen op de eeuwige ordeningen van de natuur en de menselijke natuur aan te roepen om het idee van natuurrecht te legitimeren. Het is niet alleen zo dat historische traditie en de natuurlijke orde van het leven de eenheid van de wereld vormen waarin we als mensen leven; de manier waarop we elkaar ervaren, de manier waarop we historische tradities ervaren, de manier waarop we de natuurlijke gegevenheid van ons bestaan en van onze wereld ervaren, vormen een waarlijk hermeneutisch universum, waarin we niet gevangen zitten, alsof achter onoverkomelijke barrières, maar waartoe we openstaan.

Een reflectie over wat waarheid is in de geesteswetenschappen mag niet proberen zichzelf door reflectie los te maken van de traditie waarvan ze de bindende kracht heeft erkend. Daardoor moet ze in haar eigen werk trachten zoveel mogelijk historische zelftransparantie te verwerven. In haar streven om het universum van het begrijpen beter te begrijpen dan onder het moderne wetenschappelijke kennisbegrip mogelijk lijkt, moet ze proberen een nieuwe relatie tot de begrippen die ze gebruikt te vestigen. Ze moet zich bewust zijn van het feit dat haar eigen begrip en interpretatie geen constructies zijn die op principes zijn gebaseerd, maar de voortzetting van een gebeurtenis die ver teruggaat. Daardoor zal ze de begrippen die ze gebruikt niet klakkeloos kunnen overnemen, maar zal ze de kenmerken van de oorspronkelijke betekenis van haar begrippen die tot haar zijn gekomen, moeten integreren.

De filosofische inspanning van onze tijd verschilt van de klassieke traditie van de filosofie doordat ze geen directe en ononderbroken voortzetting daarvan is. Ondanks haar verbondenheid met haar historische oorsprong is de hedendaagse filosofie zich terdege bewust van de historische afstand tussen haar en haar klassieke modellen. Dit komt met name tot uiting in haar veranderde houding ten opzichte van het begrip. Hoe belangrijk en fundamenteel de transformaties ook waren die plaatsvonden met de latinisering van Griekse begrippen en de vertaling van Latijnse begrippen in de moderne talen, de opkomst van het historische bewustzijn in de afgelopen paar eeuwen is een veel radicalere breuk. Sindsdien is de continuïteit van de westerse filosofische traditie alleen nog maar fragmentarisch effectief. We zijn die naïeve onschuld kwijtgeraakt waarmee traditionele begrippen zonder meer werden gebruikt voor het eigen denken. Sinds die tijd is de houding van de wetenschap ten opzichte van deze begrippen vreemd afstandelijk geworden, of ze nu op een geleerde, zo niet zelfbewust verouderde manier overneemt, of ze als instrumenten behandelt. Geen van beide benaderingen doet recht aan de hermeneutische ervaring. De begrippenwereld waarin het filosofisch denken zich ontwikkelt, heeft ons al op dezelfde manier in haar greep als de taal waarin we leven ons conditioneert. Als het denken verantwoordelijk wil zijn, moet het zich bewust worden van deze voorafgaande invloeden. Een nieuw kritisch bewustzijn moet nu alle verantwoordelijke filosofie begeleiden, die de denkwijzen en taalkundige gewoonten die in het individu zijn gevormd door zijn omgang met zijn omgeving, voor het forum van de historische traditie plaatst waartoe we allemaal behoren.

Het volgende onderzoek probeert aan deze eis te voldoen door een onderzoek naar de geschiedenis van begrippen zo nauw mogelijk te verbinden met de inhoudelijke uiteenzetting van haar thema. Die gewetensvolheid van fenomenologische beschrijving die Husserl voor ons allemaal tot een plicht heeft gemaakt; de breedte van de historische horizon waarin Dilthey al het filosofisch denken heeft geplaatst; en, niet in de laatste plaats, de doordringing van beide invloeden door de impuls die van Heidegger is uitgegaan, geven de maatstaf aan waarlangs de schrijver gemeten wil worden, en die hij, ondanks alle onvolkomenheden in de uitvoering, zonder voorbehoud zou willen zien toegepast.
